% generated from papers/XV-spectral/paper_XV_spectral.tex -- edit the paper, then rerun assemble.py
\chapter[The spectral route to Cartan uniqueness]{What $\sigma$-equivariance gives for free,\\ and why the range projections cannot give the rest}\label{ch:XV}
\chapterorigin{This chapter is Paper XV of the series (\texttt{papers/XV-spectral/}).}
\begin{summary}
We examine the proposal to prove Cartan uniqueness for $(\cA_{K,S},\sigma)$ by spectral means
alone, and separate what it yields from what it cannot.

On the positive side, $\sigma$-equivariance already gives more than Chapter~\ref{ch:XIII}
(its Theorem 2.5) extracted from it, and gives it without any Cartan input. The eigenvalue
group $\Sp_p(\sigma)$ of the flow is the subgroup of $\R_{>0}$ generated by
$\{\Ns\pp\}_{\pp\in S}$; and the entire $\mathrm{KMS}$
bundle is defined from $(\cA,\sigma)$ alone, so a $\sigma$-equivariant isomorphism transports
it. Consequently the function $\beta\mapsto G_S/\Xi_\beta^\perp$, up to homeomorphism, is an
invariant: in particular $|\Xi_\beta|$ wherever it is finite, the transition locus, and
$\beta_c$. \emph{No Cartan subalgebra is needed for any of this.} What Cartan uniqueness is
needed for is exactly one thing: the norm \emph{multiset}, since $\Sp_p(\sigma)$ sees only the
group generated by the norms and is blind to multiplicity --- two equal-norm primes and one
prime contribute the same subgroup.

On the negative side, the proposed route cannot deliver it, and the obstruction is at Step 3
rather than Step 2. The semilattice $\cE=\{e_\aaa\}$ of range projections does \emph{not}
generate $C(Y_S)$: since $e_\aaa=\mathbf 1_{\aaa Y_S}$ is the condition $v_\pp(x)\ge a_\pp$,
these functions depend only on the valuations and take equal values at $[x,g]$ and $[x,g']$
for all $g,g'\in G_S$. They do not separate points, and $C^*(\cE)\cong C(\overline\N^{\,S})$ is
a proper subalgebra, blind to the whole $G_S$-direction --- which by the exact sequence
$1\to U_S/\overline{\cO^\times_{K,+}}\to G_S\to\Cl^+_S\to1$ is most of $Y_S$. A complete
spectral recovery of the range projections would therefore still not give $C(Y_S)$, and this
failure is independent of whether Step 2 holds.

On Step 2 itself the verdict is negative, and no ergodic theory is needed. For equal-norm
primes $\pp_1,\pp_2$ the element $v=\mu_{\pp_1}e_{\pp_1}+\mu_{\pp_2}(1-e_{\pp_1})$, with
$e_{\pp_1}=\mu_{\pp_1}\mu_{\pp_1}^*$, is an isometry of pure frequency ---
$\mu_{\pp_1}^*(1-e_{\pp_1})=0$, so the cross terms vanish identically --- and it is not
diagonal. Mixed isometries therefore exist for every equal-norm pair. The general criterion,
for $v=\mu_{\pp_1}p+\mu_{\pp_2}(1-p)$ with $g=\pp_2\pp_1^{-1}$ and $U=\mathrm{supp}(1-p)$, is
$v^*v=1\iff g\,(U\cap\pp_1Y_S)\subseteq U$; whether the measure-preserving element $g$
($\Ns g=1$) is ergodic on $(\widetilde Y_S,\nu_\beta)$ remains open, but it no longer
governs the existence question.

The gauge-equivariant theorem of Chapter~\ref{ch:XIV} (its Theorem 4.2) therefore remains the strongest
unconditional statement, and the norm multiset remains the single quantity that
$\sigma$-equivariance alone has not been shown to recover.
\end{summary}


\section{Spectral subspaces and what they see}

We use the notation of [Ch.~\ref{ch:Main}]: $\cA_{K,S}=C(Y_S)\rtimes J_S$, spanned by the monomials
$\mu_\aaa f\mu_\bb^*$, with $\sigma_t(\mu_\aaa)=\Ns\aaa^{it}\mu_\aaa$; $\widetilde Y_S$ and
the quasi-invariant measures $\nu_\beta$ are as there.For $\lambda>0$ let
\[
\cA(\lambda)=\bigl\{x\in\cA_{K,S}:\sigma_t(x)=\lambda^{it}x\ \forall t\bigr\}
=\overline{\spn}\bigl\{\mu_\aaa f\mu_\bb^*:\Ns\aaa/\Ns\bb=\lambda,\ f\in C(Y_S)\bigr\}
\]
be the spectral subspaces of the flow in the sense of Arveson \cite{Arveson}.

\begin{proposition}[The point spectrum]\label{XV:prop:spectrum}
$\Sp_p(\sigma)=\{\lambda>0:\cA(\lambda)\ne0\}$ is the subgroup of $(\R_{>0},\times)$ generated
by $\{\Ns\pp\}_{\pp\in S}$. It is an invariant of the $C^*$-dynamical system, requiring no
Cartan subalgebra. The Arveson spectrum proper --- the closed support of the spectral
subspaces --- is the closure of this subgroup, hence all of $\R_{>0}$ as soon as two of the
$\log\Ns\pp$ have irrational ratio; the eigenvalue group $\Sp_p(\sigma)$, not its closure, is
the invariant that retains the information.
\end{proposition}

\begin{proof}
Every $\lambda=\Ns\aaa/\Ns\bb$ is realized by the eigenvector $\mu_\aaa f\mu_\bb^*$, and these
ratios form exactly the subgroup generated by the $\Ns\pp$. Conversely let
$0\ne x\in\cA(\lambda)$ and pick a functional $\varphi$ with $\varphi(x)\ne0$. Then
$t\mapsto\varphi(\sigma_t(x))=\lambda^{it}\varphi(x)$ is the uniform limit of the functions
$t\mapsto\varphi(\sigma_t(x_n))$ for finite sums $x_n$ of monomials, which are trigonometric
polynomials with frequencies in $\log\bigl\langle\Ns\pp\bigr\rangle$; the Bohr--Fourier
coefficient at $\log\lambda$ is continuous for uniform convergence, vanishes on every $x_n$
if $\lambda\notin\langle\Ns\pp\rangle$, and equals $\varphi(x)\ne0$ on the limit. So
$\lambda\in\langle\Ns\pp\rangle$. Invariance is immediate from
$\Phi\circ\sigma^1_t=\sigma^2_t\circ\Phi$.
\end{proof}

\begin{remark}[What the spectrum misses]\label{XV:rem:missing}
$\Sp_p(\sigma)$ records the \emph{group} generated by the norms, not the multiset. If
$\pp_1\ne\pp_2$ have $\Ns\pp_1=\Ns\pp_2=q$, the subgroup contributed is $q^\Z$, the same as
for a single prime of norm $q$. Multiplicity is exactly the information that the spectrum
discards, and by \S\ref{XV:sec:free} it is the only piece of the reconstruction that is not
already free.
\end{remark}

\section{What $\sigma$-equivariance gives for free}\label{XV:sec:free}

\begin{proposition}[The KMS bundle is intrinsic]\label{XV:prop:kmsintrinsic}
The set $\mathrm{KMS}_\beta(\cA,\sigma)$ is defined from $(\cA,\sigma)$ alone. Hence a
$\sigma$-equivariant isomorphism $\Phi$ induces, for every $\beta$, an affine homeomorphism
$\mathrm{KMS}_\beta(\cA_1,\sigma_1)\to\mathrm{KMS}_\beta(\cA_2,\sigma_2)$, with no hypothesis
on Cartan subalgebras.
\end{proposition}

\begin{theorem}[Free consequences]\label{XV:thm:free}
Let $\Phi:(\cA_{K_1,S_1},\sigma_1)\to(\cA_{K_2,S_2},\sigma_2)$ be a $\sigma$-equivariant
isomorphism. Then, unconditionally --- for all number fields, all prime sets, separated or
not:
\begin{enumerate}
\item $\Sp_p(\sigma_1)=\Sp_p(\sigma_2)$, so the subgroups of $\R_{>0}$ generated by the norms
agree;
\item for every $\beta>0$ the spaces $G_{S_1}/\Xi_\beta(K_1,S_1)^\perp$ and
$G_{S_2}/\Xi_\beta(K_2,S_2)^\perp$ are homeomorphic, being the extreme boundaries of the
corresponding simplices $\Prob(G_S/\Xi_\beta^\perp)$ of [Ch.~\ref{ch:Main}, Thm.~3.6]; in particular,
taking any $\beta>\beta_c$, the compact spaces $G_{S_1}$ and $G_{S_2}$ are homeomorphic;
\item in particular $|\Xi_\beta(K_1,S_1)|=|\Xi_\beta(K_2,S_2)|$ whenever these are finite;
\item the transition loci coincide, and $\beta_c(S_1)=\beta_c(S_2)$.
\end{enumerate}
\end{theorem}

\begin{proof}
(1) is Proposition~\ref{XV:prop:spectrum}. (2)--(4) follow from
Proposition~\ref{XV:prop:kmsintrinsic}: the extreme boundary of a Bauer simplex is a
homeomorphism invariant, and by [Ch.~\ref{ch:Main}, Thm.~3.6] it is $G_S/\Xi_\beta^\perp$, which for
$\beta>\beta_c$ is $G_S$ itself since $\zeta_{K,S}(\beta)<\infty$ forces
$\Xi_\beta=\widehat{G_S}$; the transition locus is the set of $\beta$ at which the simplex
changes, and $\beta_c$ is determined with it --- this is the argument of Theorem 2.5(5) of
Chapter~\ref{ch:XIII}, which uses only the transport of $\mathrm{KMS}$ states available
here and no Cartan input.
\end{proof}

\begin{remark}[The gap is exactly the multiset]\label{XV:rem:onlygap}
Comparing Theorem~\ref{XV:thm:free} with Theorem 2.5 of Chapter~\ref{ch:XIII}, the conclusions there not
obtained here are the norm-preserving bijection $S_1\to S_2$ of its parts (2)--(3) and, in
its conditional part (6), the matching of Frobenius classes and the transport of
$\{\Xi_\beta\}$ as subgroups of $\widehat{G_S}$ --- here only the cardinalities and the
quotient spaces come for free. The narrow class number is a conclusion in neither place: by
Remark 3.2 of Chapter~\ref{ch:XIII}, $h^+_S$ is recovered by no argument there, and nothing here adds
one. Everything else --- $G_S$ as a compact space, the chain $\{\Xi_\beta\}$ in cardinality,
the transition locus, $\beta_c$ --- comes free from $\sigma$-equivariance. So the Cartan
question is not a prerequisite for the arithmetic reconstruction in general; it is a
prerequisite for one specific piece of it, the multiplicity of the norms, and by
Remark~\ref{XV:rem:missing} that is precisely what $\Sp_p(\sigma)$ cannot see.
\end{remark}

\section{Mixed isometries exist: Step 2 fails}

Assume $\pp_1\ne\pp_2$ in $S$ with $\Ns\pp_1=\Ns\pp_2=q$, and put $g=\pp_2\pp_1^{-1}\in I_S$.

\begin{lemma}[Nica relation]\label{XV:lem:nica}
For coprime $\pp_1,\pp_2$ one has $\mu_{\pp_1}^*\mu_{\pp_2}=\mu_{\pp_2}\mu_{\pp_1}^*$.
\end{lemma}

\begin{proof}
$\pp_1\vee\pp_2=\pp_1\pp_2$, and the Nica covariance relation \cite{Nica}
$\mu_\aaa^*\mu_\bb=\mu_{\aaa^{-1}(\aaa\vee\bb)}\mu_{\bb^{-1}(\aaa\vee\bb)}^*$ gives the claim.
\end{proof}

\begin{proposition}[Mixed isometries exist; the general criterion]\label{XV:prop:mixed}
Let $p\in C(Y_S)$ be a projection, $U=\mathrm{supp}(1-p)$, and
$v=\mu_{\pp_1}p+\mu_{\pp_2}(1-p)$. Then
\[
v^*v=1+\mathbf 1_{U^c\cap\,g(U\cap\pp_1Y_S)}\,\mu_{\pp_2}\mu_{\pp_1}^*+\text{adjoint},
\qquad\text{so}\qquad
v^*v=1\iff g\,(U\cap\pp_1Y_S)\subseteq U,
\]
where $g$ acts as the partial homeomorphism $\pp_1Y_S\to\pp_2Y_S$. In particular the
condition is vacuous when $U\cap\pp_1Y_S=\emptyset$: for $p=e_{\pp_1}$, i.e.\
$U=Y_S\setminus\pp_1Y_S$, where $e_{\pp_1}=\mu_{\pp_1}\mu_{\pp_1}^*=\mathbf 1_{\pp_1Y_S}$,
\[
v=\mu_{\pp_1}e_{\pp_1}+\mu_{\pp_2}(1-e_{\pp_1})
\]
is an isometry, since $\mu_{\pp_1}^*(1-e_{\pp_1})=\mu_{\pp_1}^*-\mu_{\pp_1}^*\mu_{\pp_1}\mu_{\pp_1}^*=0$
kills the cross terms outright, and it is homogeneous for $\sigma$ of frequency $q$ but not
of a single gauge degree, hence not of the form $\mu_\pp f$. \emph{Mixed isometries of pure
frequency exist for every equal-norm pair, with no further hypothesis.}
\end{proposition}

\begin{proof}
Expanding and using Lemma~\ref{XV:lem:nica},
$v^*v=p+(1-p)+p\,\mu_{\pp_2}\mu_{\pp_1}^*(1-p)+\text{adjoint}$. Pushing the functions
through the partial isometry $\mu_{\pp_2}\mu_{\pp_1}^*$, which implements
$g:\pp_1Y_S\to\pp_2Y_S$, turns the third term into
$\mathbf 1_{U^c}\mathbf 1_{g(U\cap\pp_1Y_S)}\,\mu_{\pp_2}\mu_{\pp_1}^*$, which lies in the
spectral subspace of gauge degree $\pp_2-\pp_1\ne0$ and so is linearly independent of
$C(Y_S)$; it vanishes iff $U^c\cap g(U\cap\pp_1Y_S)=\emptyset$, i.e.\
$g(U\cap\pp_1Y_S)\subseteq U$. The special case is the display, and $v$ has gauge degrees
$\delta_{\pp_1}$ and $\delta_{\pp_2}$ both occurring, while any $\mu_\pp f$ is homogeneous.
\end{proof}

\begin{remark}[Status]\label{XV:rem:ergstatus}
Step 2 of the proposed route --- ``every isometry of pure frequency is diagonal'' --- is
therefore \emph{false}, by a two-line example, and the failure needs no ergodic theory: the
projection $e_{\pp_1}$, being a range projection, is triggered by the valuations alone, and
mixing over it costs nothing. An earlier draft of this note claimed that mixed isometries
exist \emph{iff} the element $g=\pp_2\pp_1^{-1}$ fails to be ergodic on
$(\widetilde Y_S,\nu_\beta)$; both implications fail --- the example above uses no invariant
set, and non-ergodicity produces measurable, not clopen, invariant sets. What survives of the
ergodic picture is this: $\Ns g=1$, so $g$ preserves each $\nu_\beta$ (the Radon--Nikodym
cocycle of the $I_S$-action being $(\Ns g)^{-\beta}$), the full group $I_S$ acts ergodically
by the Kolmogorov zero--one law [Ch.~\ref{ch:Main}, Lem.~3.1], and whether a \emph{single} such $g$
--- shifting $v_{\pp_2}$ up, $v_{\pp_1}$ down, and translating $G_S$ by
$\sigma_{\pp_2}\sigma_{\pp_1}^{-1}$ --- is ergodic remains open and of independent interest;
but it no longer decides the existence of mixed isometries, which is settled affirmatively.
\end{remark}

\section{The decisive gap: range projections are blind to $G_S$}

\begin{theorem}[$\cE$ does not generate the diagonal]\label{XV:thm:blind}
Let $\cE=\{e_\aaa=\mathbf 1_{\aaa Y_S}:\aaa\in J_S\}$. Then $\cE$ does not separate the points
of $Y_S$: for every $\aaa$ and all $g,g'\in G_S$,
\[
e_\aaa\bigl([x,g]\bigr)=e_\aaa\bigl([x,g']\bigr).
\]
Consequently $C^*(\cE)\cong C\bigl(\overline\N^{\,S}\bigr)$, the algebra of continuous
functions of the valuations alone, and is a proper subalgebra of $C(Y_S)$ whenever $G_S$ is
nontrivial.
\end{theorem}

\begin{proof}
$\aaa Y_S=\{[x,g]:v_\pp(x)\ge a_\pp\ \forall\pp\}$, a condition on $x$ only, so $e_\aaa$
factors through the projection $Y_S\to\overline\N^{\,S}$, $[x,g]\mapsto(v_\pp(x))_\pp$. The
$e_\aaa$ generate the clopen algebra of $\overline\N^{\,S}$, whence the identification by
Stone--Weierstrass; properness follows since $G_S\ne\{1\}$ for every $S\ne\emptyset$, by
[Ch.~\ref{ch:Main}, (2.1)].
\end{proof}

\begin{corollary}[Step 3 fails independently of Step 2]\label{XV:cor:step3}
Even if every isometry of pure frequency $q$ were diagonal --- and by
Proposition~\ref{XV:prop:mixed} it is not --- so that the range projections $e_\pp$ were
spectrally recoverable, the algebra they generate is
$C(\overline\N^{\,S})$ and not $C(Y_S)$. The proposed route therefore cannot establish
$\Phi(C(Y_{S_1}))=C(Y_{S_2})$, the identity needed to bring Renault's reconstruction
\cite{Renault} to bear, and this failure is located at Step 3, independently of Step 2.
\end{corollary}

\begin{remark}[How much is missed]\label{XV:rem:howmuch}
By [Ch.~\ref{ch:Main}, (2.1)], $G_S$ sits in
$1\to U_S/\overline{\cO^\times_{K,+}}\to G_S\to\Cl^+_S\to1$ with $U_S=\prod_\pp\cO_\pp^\times$
an infinite product of infinite profinite groups. The invisible direction is therefore not a
small correction but a factor comparable in size to $Y_S$ itself. Any route to $C(Y_S)$ must
produce elements that vary along $G_S$, and the range projections, being defined by valuation
inequalities, never do.
\end{remark}

\begin{remark}[What would have to replace it]\label{XV:rem:replace}
Elements of $\cA$ that vary along $G_S$ do exist in abundance --- all of $C(Y_S)$ does ---
but they are not distinguished by the flow, $\sigma$ acting trivially on $C(Y_S)$. To
characterize them one would need a construction using the multiplication of $\cA$ against the
spectral subspaces, for instance identifying $C(Y_S)$ inside $\cA^\sigma$ as a distinguished
maximal abelian subalgebra. That is the question left open in Remark 3.3 of Chapter~\ref{ch:XIV}, and
the present note does not advance it; what it does is show that the range projections are not
a route to it.
\end{remark}

\section{Assessment}

\[
\begin{array}{lll}
\text{Statement} & \text{Status} & \text{Reference}\\\hline
\Sp_p(\sigma)=\langle\Ns\pp\rangle\le\R_{>0} & \text{true, intrinsic} & \text{Prop.~1.1}\\
\Sp_p(\sigma)\ \text{sees the norm multiset} & \text{false: no multiplicity} & \text{Rem.~1.2}\\
\text{KMS bundle intrinsic to }(\cA,\sigma) & \textbf{true} & \text{Prop.~2.1}\\
G_S,\ \{|\Xi_\beta|\},\ \text{transition loci},\ \beta_c\ \text{recovered} & \textbf{true, unconditionally} & \text{Thm.~2.2}\\
h^+_S\ \text{recovered} & \text{not claimed (as in Chapter~\ref{ch:XIII})} & \text{Rem.~2.3}\\
\text{norm multiset recovered} & \textbf{open; the only gap} & \text{Rem.~2.3}\\
\text{mixed isometries exist} & \textbf{true, unconditionally} & \text{Prop.~3.2}\\
\text{every pure-frequency isometry is diagonal} & \textbf{false} & \text{Prop.~3.2}\\
\text{mixed isometries}\iff g\ \text{non-ergodic} & \text{false, both ways} & \text{Rem.~3.3}\\
\cE\ \text{generates}\ C(Y_S) & \textbf{false: blind to }G_S & \text{Thm.~4.1}\\
\text{spectral route to Cartan uniqueness} & \textbf{closed} & \text{Cor.~4.2}\\
\sigma\text{-equivariance}\Rightarrow\text{gauge-equivariance} & \text{not established} & \text{---}
\end{array}
\]

Two conclusions, of opposite sign, and the positive one is the more useful.

The reconstruction of Chapter~\ref{ch:XIII} (its Theorem 2.5) was presented as resting on Cartan
uniqueness throughout. Theorem~\ref{XV:thm:free} shows that most of it does not: the
$\mathrm{KMS}$ bundle is manifestly intrinsic to $(\cA,\sigma)$, and with it come $G_S$ as a
compact space, the cardinalities of the obstruction groups, the transition locus and the
critical exponent, for every number field and every prime set, separated or not. The Cartan
question is needed for one thing only, the multiplicity of the norms, which is exactly what
the point spectrum discards.

The proposed spectral route to that last piece is closed twice over. Step 2 is false
outright: mixed isometries of pure frequency exist for every equal-norm pair
(Proposition~\ref{XV:prop:mixed}). And even its truth would have delivered only the semilattice
$\cE$, which by Theorem~\ref{XV:thm:blind} is blind to $G_S$, generating a proper subalgebra of
the diagonal. The unconditional statement therefore remains Theorem 4.2 of Chapter~\ref{ch:XIV}, with
gauge-equivariance as hypothesis.

Two questions. Is $g=\pp_2\pp_1^{-1}$ ergodic on $(\widetilde Y_S,\nu_\beta)$? By
Remark~\ref{XV:rem:ergstatus} this no longer decides the existence of mixed isometries, but it
is a clean self-contained problem about a single measure-preserving transformation of a
product of geometric measures twisted by a group translation, and we have not decided it.
And is there any construction recovering
the $G_S$-direction of $C(Y_S)$ from $(\cA,\sigma)$? Remark~\ref{XV:rem:replace} explains why the
obvious candidates do not, the flow being trivial on precisely the algebra one is trying to
find.

